\documentclass[11pt,a4paper]{article}

\usepackage[margin=0.82in]{geometry}
\usepackage[T1]{fontenc}
\usepackage{lmodern}
\usepackage{microtype}
\usepackage{parskip}
\usepackage{titlesec}

\titleformat{\section}{\large\bfseries}{}{0pt}{}
\titlespacing*{\section}{0pt}{11pt}{5pt}
\pagestyle{empty}

\begin{document}

{\LARGE Research Synopsis\par}
\vspace{7pt}
{\large Physical-Only Proprioceptive Recognition of Fifteen Classes with a LEAP Hand\par}
\vspace{12pt}
\noindent\textbf{Category}\quad Robotics and Intelligent Machines (ROBO)\par
\noindent\textbf{Student researcher}\quad G. Sai Pranav\par
\noindent\textbf{Guide}\quad Harish J. Palanthandalam-Madapusi\par

\section*{Project Abstract}

People can often distinguish objects by how their fingers close around them. This study asks whether a robotic hand can make a limited version of that judgment using its own joint positions, without camera images or dedicated tactile arrays. A physical LEAP Hand with four fingers and 16 actuated joints performed the same sequence of three controlled grips on 14 household objects and an empty-hand condition. Motor-current thresholds stopped fingers on resistance, but recognition used only the 16 final joint angles from each grip, giving 48 measurements per trial. Five physical trials per label produced 75 trials across 15 labels.

A small neural classifier was evaluated by three-fold within-dataset validation. Its highest-probability label was correct for 70 of 75 held-out trials (93.3\%). A decision rule accepted 67 trials as known classes, of which 65 were correct (97.0\% of accepted decisions); eight were uncertain. A nearest-trial method classified 69 of 75 correctly in a separate leave-one-trial-out evaluation. An exploratory 20-label extension, which contains the same 75 trials plus 25 more, gave 85 of 100 correct raw neural predictions. A further exploratory nearest-trial check found 72 of 75 correct using only the first grip, so the benefit of combining three grips is not established by these data.

The findings show that repeatable physical grasp positions contain useful information about the enrolled items. The 15-label set was selected after preliminary analysis, and the validation folds also guided model training and decision thresholds. These are internal, post-selection estimates; performance across later sessions, new object instances, and genuinely unseen classes remains unmeasured.

\section*{Introduction and Objective}

A dexterous robotic hand may learn about an object through the positions its joints reach when the fingers close around it. Prior studies have used joint information to infer object shape or identity, yet the value of a simple, repeatable sequence of final grasp angles on a physical LEAP Hand remains worth testing. This project asks whether those angles can distinguish common household items, including an empty-hand condition, without supplying images or tactile-array data to the recognition model. The objective was to collect physical trials under a consistent grip protocol, train an angle-only classifier, and compare its internal validation result with a nearest-trial method. The study also examined when a prediction should be withheld as uncertain. Generalization to new sessions and new physical exemplars was left for later work.

\section*{Innovation}

The contribution is a physical-hand recognition pipeline that turns three controlled, sequential grasps into a compact 48-angle object signature. Motor current stops each finger when resistance is detected, while only final joint angles reach the classifier. This keeps contact control separate from recognition input. The project combines a neural classifier, a nearest-trial comparison, an uncertainty decision rule, and resumable physical-data collection. The novelty lies in this specific implemented protocol and evaluation, rather than in a claim that proprioceptive object recognition itself is new.

\newpage

\section*{Methodology}

A four-finger LEAP Hand with 16 actuated joints ran the same three-grip program for every trial. Each item was placed near a consistent point relative to the palm. During each grip, joints advanced toward configured targets. When a moving joint's signed motor current exceeded its threshold, all four joints of that finger stopped; the other fingers continued until resistance or an angle target stopped them. The final 16 joint angles were read after each grip and concatenated in grip order. Motor currents were logged but excluded from the recognition features.

The main dataset comprised five physical trials for each of 14 household objects plus an empty-hand class, giving 75 labeled 48-angle vectors. Each angle was normalized using its grip-specific start and target values. The neural classifier had 48 inputs, one 24-unit tanh hidden layer, and 15 softmax outputs. Training used full-batch Adam; noisy angle copies were added only to training folds. Three-fold validation held out real physical trials from every class. A calibrated decision rule returned a class only when confidence and distance checks passed, otherwise marking it uncertain or unknown. A nearest-trial method provided a separate leave-one-trial-out comparison. We also examined a 20-label extension containing the same 75 trials and 25 more, and recalculated nearest-trial accuracy for subsets of the three grips. The selected 15-label set was not fixed before preliminary analysis.

\section*{Results and Conclusions}

The neural model's top label was correct for 70 of 75 held-out physical trials (93.3\%). It accepted 67 as known classes, correctly labeling 65 of them (97.0\% of accepted decisions); eight were uncertain. The nearest-trial method correctly labeled 69 of 75 trials (92.0\%) in its different leave-one-trial-out evaluation. In the 20-label extension, raw neural accuracy was 85/100 and 61/64 accepted decisions were correct. Exploratory nearest-trial reanalysis found 72/75 correct using grip 1 alone versus 69/75 using all three, so an accuracy benefit from more grips was not demonstrated. These results show useful class information in final angles, but class selection, model settings and decision thresholds were informed by the same data. New sessions, new physical object instances, unseen classes and repeat full sequences require separate tests.

\section*{Acknowledgement and References}

I thank Harish J. Palanthandalam-Madapusi, the project's single guide, for research guidance and technical feedback. Key references are Shaw, Agarwal and Pathak (2023), LEAP Hand, arXiv:2309.06440; V\'asquez and Perdereau (2017), \emph{Robotics and Autonomous Systems} 98, 135--146; Sintov and Meir (2023), \emph{International Journal of Robotics Research} 42, 537--561; and Kutsuzawa et al.\ (2024), \emph{Frontiers in Neurorobotics} 18, 1466630. The complete references and technical details appear in the research paper.

\end{document}
